% BHU PYQ -- Manually transcribed & error-corrected
% Paper: BCF-116 : Business Statistics
% Exam: B.Com. (Hons.) FMM Semester I Examination, 2013-14

\documentclass[11pt,a4paper]{article}
\usepackage[a4paper,top=2.5cm,bottom=2.5cm,left=2.8cm,right=2.8cm,headheight=14pt]{geometry}
\usepackage{amsmath,amssymb}
\usepackage[T1]{fontenc}
\usepackage[utf8]{inputenc}
\usepackage{lmodern,microtype,fancyhdr,enumitem,hyperref,lastpage,parskip}

\hypersetup{
    colorlinks=true,
    urlcolor=black,
    linkcolor=black,
    pdftitle={BCF-116: Business Statistics -- BHU B.Com. (Hons.) FMM Sem I 2013-14}
}

\pagestyle{fancy}
\fancyhf{}
\renewcommand{\headrulewidth}{0.4pt}
\renewcommand{\footrulewidth}{0.4pt}
\fancyhead[L]{\small   \textit{Banaras Hindu University}}
\fancyhead[R]{\small   \textit{BCF-116: Business Statistics}}
\fancyfoot[C]{\small Page  \thepage\ of \pageref{LastPage}}


\newlist{parts}{enumerate}{2}
\setlist[parts,1]{label=(\alph*),leftmargin=*,itemsep=5pt,topsep=3pt}
\setlist[parts,2]{label=(\roman*),leftmargin=*,itemsep=3pt,topsep=2pt}

\newcommand{\pts}[1]{\hfill   \textit{[#1]}}


\begin{document}


\begin{center}
  {\large\bfseries BANARAS HINDU UNIVERSITY}\\[2pt]
  {
\normalsize B.Com. (Hons.) FMM --- Semester I Examination, 2013-14}\\[6pt]
  \rule{0.8\linewidth}{0.5pt}\\[6pt]
  {\large\bfseries Paper No. BCF-116: Business Statistics}\\[4pt]
  {
\normalsize Time: 3 Hours\hfill Maximum Marks: 70}
\end{center}

\rule{\linewidth}{0.4pt}

\smallskip



\noindent   \textit{Instructions: Answer all questions. All questions carry equal marks (14 marks each).}

\bigskip




\noindent   \textbf{Question 1.}
``Data are like clay out of which you can make a God or a Devil as you please.'' Discuss.\pts{14}

\centerline{   \textbf{OR}}

Discuss in brief, the methods generally used in the collection of primary data.\pts{14}

\medskip\rule{\linewidth}{0.2pt}

\pagebreak




\noindent   \textbf{Question 2.}
An incomplete distribution is given as follows:


\begin{center}
\small

\begin{tabular}{|c|c|}
\hline
   \textbf{Class Interval} &   \textbf{Frequency} \\
\hline
10--20 & 12 \\
20--30 & 30 \\
30--40 & ? \\
40--50 & 65 \\
50--60 & ? \\
60--70 & 25 \\
70--80 & 18 \\
\hline
   \textbf{Total ($\sum f$)} &   \textbf{229} \\
\hline
\end{tabular}
\end{center}

Find the missing frequencies, if the median is 46.\pts{14}

\centerline{   \textbf{OR}}

What is statistical average? Explain the properties of a good average.\pts{14}

\medskip\rule{\linewidth}{0.2pt}




\noindent   \textbf{Question 3.}
What do you mean by dispersion? Discuss the various methods of measuring dispersion. Also discuss their relative merits and demerits.\pts{14}

\centerline{   \textbf{OR}}

Calculate standard deviation from the following series:


\begin{center}
\small

\begin{tabular}{|l|c|c|c|c|c|c|c|c|}
\hline
   \textbf{Marks (more than)} & 0 & 10 & 20 & 30 & 40 & 50 & 60 & 70 \\
\hline
   \textbf{No. of students} & 100 & 90 & 75 & 50 & 20 & 10 & 5 & 0 \\
\hline
\end{tabular}
\end{center}\pts{14}

\medskip\rule{\linewidth}{0.2pt}

\pagebreak




\noindent   \textbf{Question 4.}
What is `time series'? Describe briefly the various methods of measurement of secular trend.\pts{14}

\centerline{   \textbf{OR}}

Fit a straight line trend by the method of least squares in the following series:


\begin{center}
\small

\begin{tabular}{|l|c|c|c|c|c|c|}
\hline
   \textbf{Year} & 2001 & 2002 & 2003 & 2004 & 2005 & 2006 \\
\hline
   \textbf{Production} & 7 & 10 & 12 & 14 & 17 & 24 \\
\hline
\end{tabular}
\end{center}\pts{14}

\medskip\rule{\linewidth}{0.2pt}




\noindent   \textbf{Question 5.}
Define index number and describe in detail its uses.\pts{14}

\centerline{   \textbf{OR}}

From the data given below, construct Fisher's ideal index number:


\begin{center}
\small

\begin{tabular}{|c|c|c||c|c|}
\hline
   \textbf{Commodities} & \multicolumn{2}{c||}{   \textbf{2009}} & \multicolumn{2}{c|}{   \textbf{2010}} \\
\cline{2-5}
 &   \textbf{Price (per unit)} &   \textbf{Expenditure (in Rs.)} &   \textbf{Price (per unit)} &   \textbf{Expenditure (in Rs.)} \\
\hline
A & 2 & 40 & 5 & 75 \\
B & 4 & 16 & 8 & 40 \\
C & 1 & 10 & 2 & 24 \\
D & 5 & 25 & 10 & 60 \\
\hline
\end{tabular}
\end{center}\pts{14}

\vfill
\rule{\linewidth}{0.4pt}

\begin{center}\small
\href{https://bhu-exam.vercel.app/}{Click here to visit our website}\\[4pt]
\href{https://me-aryan.vercel.app/}{Click here to visit our contributor}\\[4pt]
\href{https://quashan-me.vercel.app/}{Click here to visit our contributor}
\end{center}

\end{document}
